\subsection{Funções Acadêmicas e Pesagem (C012, C016)}

Como as regras do Firestore bloqueiam a escrita direta nos nós de Turmas, Comentários e Roteiros (\texttt{allow write: if false;}), as operações de criação, moderação e compartilhamento são intermediadas por \textit{Cloud Functions} para garantir trilha de auditoria e sanitização de entradas.

\begin{lstlisting}[language=TypeScript, title={Moderação de Comentários e Posts}]
export const adicionarComentario = onCall(async (request) => {
  const { idTurma, idPost, texto } = request.data as { idTurma: string; idPost: string; texto: string };
  await validarVinculoTurma(request, idTurma); // Lança erro se não matriculado/professor

  const comentarioRef = admin.firestore().collection(`Turma/${idTurma}/Posts/${idPost}/Comentarios`).doc();
  await comentarioRef.set({
    id_autor: request.auth!.uid,
    texto: texto,
    moderado: false,
    criado_em: admin.firestore.FieldValue.serverTimestamp(),
  });
  return { id: comentarioRef.id };
});

export const moderarComentario = onCall(async (request) => {
  const { idTurma, idPost, idComentario, motivo } = request.data;
  await validarPermissaoAcademica(request, idTurma); // Professor ou Chefe Geral

  const comentarioRef = admin.firestore().collection(`Turma/${idTurma}/Posts/${idPost}/Comentarios`).doc(idComentario);
  await admin.firestore().runTransaction(async (tx) => {
    const snap = await tx.get(comentarioRef);
    if (!snap.exists) throw new HttpsError("not-found", "Comentário inexistente.");
    
    // Grava histórico de moderação
    tx.set(comentarioRef.collection("Historico").doc(), {
      texto_original: snap.data()!.texto,
      motivo: motivo,
      moderado_por: request.auth!.uid,
      data: admin.firestore.FieldValue.serverTimestamp()
    });

    // Mascara o conteúdo atual
    tx.update(comentarioRef, {
      texto: "[Comentário removido pela moderação]",
      moderado: true,
      data_moderacao: admin.firestore.FieldValue.serverTimestamp()
    });
  });
});
\end{lstlisting}

\begin{lstlisting}[language=TypeScript, title={Compartilhamento de Roteiros}]
export const compartilharRoteiro = onCall(async (request) => {
  const { idRoteiro, idTurma } = request.data;
  await validarPermissaoAcademica(request, idTurma); // Somente professor da turma
  
  const roteiroRef = admin.firestore().collection("Roteiro").doc(idRoteiro);
  await admin.firestore().runTransaction(async (tx) => {
    const snap = await tx.get(roteiroRef);
    if (snap.data()!.id_autor !== request.auth!.uid && !request.auth!.token.Chefe_Geral) {
      throw new HttpsError("permission-denied", "Somente o autor pode compartilhar.");
    }
    tx.set(admin.firestore().collection(`Turma/${idTurma}/Roteiros_Compartilhados`).doc(idRoteiro), {
      adicionado_por: request.auth!.uid,
      data: admin.firestore.FieldValue.serverTimestamp()
    });
  });
});
\end{lstlisting}

\begin{lstlisting}[language=TypeScript, title={Pesagem de Rotina (UI-06)}]
export const registrarPesagemRotina = onCall(async (request) => {
  const { idFrasco, pesoAtual, motivo } = request.data as { idFrasco: string; pesoAtual: number; motivo: string };
  await validarPermissao(request, ["Gestor_Almoxarifado", "Chefe_Geral"]);

  const frascoRef = admin.firestore().collection("Frasco_Reagente").doc(idFrasco);
  await admin.firestore().runTransaction(async (tx) => {
    const snap = await tx.get(frascoRef);
    if (!snap.exists) throw new HttpsError("not-found", "Frasco não encontrado.");
    
    tx.update(frascoRef, {
      peso_atual: pesoAtual,
      ultima_pesagem: admin.firestore.FieldValue.serverTimestamp()
    });
    
    tx.set(admin.firestore().collection("Historico_Frasco_Reagente").doc(), {
      id_frasco_reagente: idFrasco,
      tipo: "PESAGEM_ROTINA",
      peso_registrado: pesoAtual,
      motivo: motivo,
      id_gestor: request.auth!.uid,
      timestamp: admin.firestore.FieldValue.serverTimestamp()
    });
  });
});
\end{lstlisting}
